\section{Leading terms through the entire divisibility tower}\label{sec:leading}

All degrees in this section are total polynomial degrees. Saying that
$P$ has unique highest homogeneous term $cy^D$ means that
$P-cy^D$ has degree at most $D-1$. This stronger assertion, not merely
the value of $\deg P$, prevents hidden cancellations after substitution.

\begin{lemma}\label{lem:difference}
Let $P\in R$ have unique highest homogeneous term $cy^D$, where
$c\ne0$ and $p\nmid D$. Then $\delta(P)$ has unique highest homogeneous
term
\begin{equation}\label{eq:one-difference}
 c\overline D b\,y^{q(D-1)+m}.
\end{equation}
Here $\overline D$ is the residue of $D$ in $\F_p$.
\end{lemma}

\begin{proof}
Write $P=cy^D+P_0$ with $\deg P_0\le D-1$, and set $h=g(y)-ax$.
The coordinate degrees of $H$ and $\Phi$ are at most $q$, so both
$T(P_0)$ and $U(P_0)$ have degree at most $q(D-1)$. The contribution
of the top monomial is
\[
 c\bigl((y^q+h)^D-y^{qD}\bigr)
 =c\sum_{j=1}^{D}\binom Dj y^{q(D-j)}h^j.
\]
Since $\deg h=m$ and its unique top term is $by^m$, the term with
$j=1$ has the nonzero unique top monomial in \eqref{eq:one-difference}.
For $j\ge2$, the degree is at most
\[
 q(D-j)+jm=q(D-1)+m-(j-1)(q-m),
\]
strictly smaller because $m<q$. The contribution from $P_0$ also has
strictly smaller degree. Thus no other term can cancel the stated
monomial.
\end{proof}

\begin{lemma}\label{lem:delta-iterates}
For every $j\ge1$, the unique highest homogeneous term of $\delta^j(y)$
is $c_jy^{D_j}$, where
\begin{align}
 c_j&=b^j\overline m^{\,j-1}\ne0,\label{eq:coefficient}\\
 D_1&=m,\qquad D_{j+1}=qD_j-(q-m),\label{eq:degree-recurrence}\\
 D_j&=\frac{(m-1)q^j+(q-m)}{q-1}.
 \label{eq:degree-closed}
\end{align}
In particular, $D_j\equiv m\pmod p$ and $0<D_j<q^j$.
\end{lemma}

\begin{proof}
For $j=1$ the assertion follows from \eqref{eq:delta-basic}, since
$m\ge2$. If the assertion holds at $j$, then
$D_j\equiv m\pmod p$, so Lemma~\ref{lem:difference} applies. It gives
\[
 D_{j+1}=q(D_j-1)+m,\qquad
 c_{j+1}=c_j\overline m b.
\]
The degree recurrence again gives $D_{j+1}\equiv m\pmod p$, closing
the induction. The coefficient recurrence yields \eqref{eq:coefficient}.
Solving the degree recurrence gives
\[
 D_j=q^j-(q-m)(1+q+\cdots+q^{j-1}),
\]
equivalent to \eqref{eq:degree-closed}. Positivity and the strict upper
bound follow either from this expression and \eqref{eq:degree-closed},
or by induction using $1<m<q$.
\end{proof}

The hypothesis $p\nmid m$ is used at every step, not only at the first
period: it ensures that the first noncancelled binomial term always
has nonzero coefficient. Its failure is treated explicitly in
Example~\ref{ex:failure}.

\begin{proposition}\label{prop:period-degree}
For $n=rs$, where $r=p^{v_p(n)}$ and $p\nmid s$, the polynomial
$T^n(y)-U^n(y)$ has unique highest homogeneous term
\begin{equation}\label{eq:period-term}
 s c_r y^{q^{n-r}D_r}.
\end{equation}
The polynomial $T^n(x)-U^n(x)$ has unique highest homogeneous term
$-x^{q^n}$.
\end{proposition}

\begin{proof}
Commutativity of $T$ and $U$ and the characteristic-$p$ binomial
identity give
\[
 T^r-U^r=(T-U)^r=\delta^r.
\]
The difference-of-powers factorization in the same commuting operator
algebra then gives
\begin{equation}\label{eq:factorization}
 T^n-U^n
 =\sum_{i=0}^{s-1}T^{ri}U^{r(s-1-i)}\delta^r.
\end{equation}
For any polynomial in $R$ with unique top term $cy^D$, applying either
$T$ or $U$ gives unique top term $cy^{qD}$. For $T$, the monic $y^q$
term in the second coordinate supplies it, while every lower term
has substituted degree at most $q(D-1)$. For $U$ this is immediate.
Coefficients in $\F_q$ remain unchanged.

Lemma~\ref{lem:delta-iterates} therefore shows that every summand of
\eqref{eq:factorization}, applied to $y$, has unique highest homogeneous
term
\[
 c_r y^{q^{r(s-1)}D_r}=c_r y^{q^{n-r}D_r}.
\]
The $s$ such terms add to a nonzero coefficient, because $p\nmid s$.
This proves \eqref{eq:period-term}, including $s=1$.

Finally, $T^n(x)=T^{n-1}(y)$ has degree $q^{n-1}$ and unique top
term $y^{q^{n-1}}$, by the same substitution observation starting
with $y$. Since $U^n(x)=x^{q^n}$ has strictly larger degree, the first
equation has the asserted top term.
\end{proof}
